\documentclass[11pt,a4paper]{article}

\usepackage[utf8]{inputenc}
\usepackage[T1]{fontenc}
\usepackage[italian]{babel}
\usepackage{geometry}
\usepackage{array}
\usepackage{xcolor}
\usepackage{helvet}
\usepackage{tabularx}

\renewcommand{\familydefault}{\sfdefault}
\geometry{margin=2.1cm}
\pagestyle{empty}
\setlength{\parindent}{0pt}
\renewcommand{\arraystretch}{1.25}

\definecolor{maincolor}{HTML}{0F766E}
\definecolor{softcolor}{HTML}{EAF7F5}
\definecolor{linecolor}{HTML}{B8DCD7}

\newcommand{\infobox}[2]{%
    \textcolor{maincolor}{\bfseries #1}\\[-0.05cm]
    #2
}

\begin{document}

\colorbox{maincolor}{%
    \parbox{\dimexpr\textwidth-2\fboxsep\relax}{%
        \vspace{0.35cm}
        \centering
        {\color{white}\Large\bfseries Dichiarazione progetto}\\[0.15cm]
        {\color{white}\normalsize Tecnologie Web}
        \vspace{0.35cm}
    }%
}

\vspace{0.65cm}

\begin{tabularx}{\textwidth}{>{\bfseries\color{maincolor}}p{3.6cm} X}
Studente & Carmine Sgariglia \\
Matricola & N86 005069 \\
Nome progetto & REGEXRIDDLE \\
Traccia scelta & Piattaforma web dedicata a enigmi testuali basati su espressioni regolari, con interazione tra interfaccia utente, API e database. \\
\end{tabularx}

\vspace{0.55cm}

\colorbox{softcolor}{%
    \parbox{\dimexpr\textwidth-2\fboxsep\relax}{%
        \vspace{0.25cm}
        \infobox{Back-end}{Django e Django REST Framework sono stati utilizzati per realizzare le API REST, gestire la logica applicativa e collegare il progetto al database PostgreSQL. Simple JWT supporta l'autenticazione, drf-spectacular documenta le API e Gunicorn consente l'esecuzione nello stack di produzione.}
        \vspace{0.35cm}

        \infobox{Front-end}{L'interfaccia e' stata sviluppata con React e TypeScript. Vite gestisce ambiente di sviluppo e build, Tailwind CSS definisce lo stile, React Router organizza la navigazione, mentre Axios e TanStack Query curano lo scambio dati con il back-end.}
        \vspace{0.35cm}

        \infobox{Ambiente e distribuzione}{Docker Compose permette di avviare in modo riproducibile servizi applicativi, database e componenti di supporto. Nello stack orientato alla produzione nginx agisce da reverse proxy e serve l'applicazione front-end compilata.}
        \vspace{0.25cm}
    }%
}

\vfill

\begin{center}
    \begin{minipage}{0.88\textwidth}
        \centering
        \small
        REGEXRIDDLE separa chiaramente responsabilita' front-end e back-end: il front-end offre
        l'esperienza utente per risolvere gli enigmi, mentre il back-end espone servizi REST,
        conserva i dati e valida le principali operazioni applicative.
    \end{minipage}
\end{center}

\end{document}
